\section{Pendahuluan}
\label{sec:pendahuluan}

\subsection{Latar Belakang}
Tuliskan latar belakang topik proyek robotika atau sistem cerdas yang Anda kerjakan di sini. Latar belakang harus mencakup penjelasan mengenai permasalahan riil yang melandasi pentingnya pembuatan sistem robot ini, penelitian terdahulu yang relevan, serta solusi cerdas yang diajukan dalam proyek ini. 

Contoh pengacuan pustaka menggunakan LaTeX: Robot ini dirancang untuk mengatasi permasalahan navigasi otonom menggunakan algoritma tertentu \parencite{author2026example}.

\subsection{Rumusan Masalah}
Berdasarkan latar belakang di atas, rumusan masalah dalam penelitian/proyek ini adalah:
\begin{enumerate}
    \item Bagaimana merancang arsitektur perangkat keras dan lunak pada robot agar mampu beroperasi secara cerdas?
    \item Bagaimana performa algoritma kontrol/navigasi yang diimplementasikan pada sistem robot tersebut?
    \item Sejauh mana tingkat akurasi pembacaan sensor dalam mengenali kondisi lingkungan sekitar?
\end{enumerate}

\subsection{Tujuan Penelitian}
Tujuan dari pelaksanaan proyek UAS Intelligent Systems and Robotics ini adalah:
\begin{enumerate}
    \item Merancang dan mengimplementasikan sistem robotika terintegrasi yang mampu melakukan [tuliskan tugas spesifik robot].
    \item Menganalisis kinerja sensor dan aktuator yang digunakan pada platform robot.
    \item Menguji keandalan sistem kontrol cerdas dalam skenario lingkungan nyata.
\end{enumerate}

\subsection{Manfaat Penelitian}
Hasil dari rancang bangun sistem robotika ini diharapkan dapat memberikan manfaat sebagai berikut:
\begin{itemize}
    \item \textbf{Bagi Mahasiswa:} Meningkatkan pemahaman praktis mengenai integrasi sensor, aktuator, dan algoritma kecerdasan buatan pada sistem fisik.
    \item \textbf{Bagi Akademisi:} Menjadi referensi tambahan dalam pengembangan sistem navigasi robotika otonom skala laboratorium.
    \item \textbf{Bagi Praktisi:} Memberikan inspirasi desain sistem robotika berbiaya rendah dengan performa yang andal.
\end{itemize}
